\documentclass{article}
\usepackage{geometry}
%\geometry{hmargin={2cm,2cm}}
\usepackage{graphicx}
\usepackage{xcolor}
\usepackage{hyperref}
\definecolor{logo}{RGB}{20,99,163}
\usepackage[italian]{babel}

\begin{document}
\begin{titlepage}
    \begin{center}
        \includegraphics{../../../../Images/logo}\\
        \vspace{20px}
        \textcolor{logo}{\hrulefill}\\
        \vspace{20px}
        \textbf{\huge\textcolor{logo}{Verbale Interno del 2021/03/09}}\\
        \vspace{10px}
        \textcolor{logo}{\hrulefill}\\
        \vspace{20px}
        \href{mailto:techsweave@gmail.com}{techsweave@gmail.com}\\
        \vspace{40px}
        \textbf{Informazioni sul documento}\\
        \vspace{20px}
        \begin{tabular}{p{100px} | p{100px}}
            \textbf{Versione} & 0.0.2\\
            \textbf{Responsabile} & Manuel Varo\\
            \textbf{Redattore} & Samuele De Simone\\
            \textbf{Verificatore} & Nicolò Giaccone\\
            \textbf{Destinatari} & Prof. Tullio Vardanega \newline Prof. Riccardo Cardin, \newline TechSWEave\\
            \textbf{Stato} & Approvato\\
            \textbf{Uso} & Interno\\
        \end{tabular}
            \vspace{1ex}

    \end{center}
    \end{titlepage}
    \tableofcontents
    \section{Informazioni riunione}
    \begin{itemize}
        \item Data e ora: 09/03/2021 ore 21.30-22.45.
        \item Luogo: Canale Discord del gruppo tramite videochiamata.
        \item Partecipanti interni: Samuele De Simone, Manuel Varo, Tito Scutari, Nicolò Giaccone, Amedeo Meggiolaro, Simone Urbani, Marco Barbaresco.
        \item Partecipanti esterni: nessuno.
        \item Segretario: Amedeo Meggiolaro.
    \end{itemize}
    \section{Argomenti della riunione}
    \begin{itemize}
        \item Scelta del nome, del logo e della e-mail del gruppo.
        \item Scelta degli strumenti di comunicazione.
        \item Decisioni sulle prossime riunioni.
        \item Studio dei capitolati.
        \item Analisi dei documenti da produrre per la RR.
        \item Scelta dei ruoli.
    \end{itemize}
    \section{Resoconto}
    
        \subsection{Scelta del nome, del logo e della e-mail del gruppo}
        Sono stati proposti vari nomi con cui identificare il gruppo, dopo una breve
        discussione è stato scelto il nome TechSWEave.\\
        Sulla base di questo nome abbiamo creato la casella di posta elettronica 
        techsweave@gmail.com e Tito Scutari è stato incaricato per la creazione del logo.
        
        \subsection{Scelta degli strumenti}
        Il gruppo ha scelto i vari strumenti per gestire, comunicare e coordinarsi sul lavoro.
        Le scelte sono state fatte in base alla comodità e facilità di utilizzo dei
        vari ambienti.
        \begin{itemize}
            \item Discord: utile per tutte le riunioni in videochiamata e per parlare giornalmente
            \item Telegram: semplice chat in cui si discutono problemi secondari.
            \item GitHub: per il versionamento e la condivisione di documenti e codice.
            \item Visual Studio Code: sia per la scrittura dei pdf tramite LaTeX che per la scrittura del codice.
            \item LaTeX: per la creazione dei documenti.
        \end{itemize}

        \subsection{Decisioni sulle prossime riunioni}
        Il gruppo ha concordato di fare due riunioni per settimana in giorni prefissati.

        \subsection{Studio dei capitolati}
        Sono stati studiati i vari capitolati e ciascun membro ha scelto di analizzarne uno in particolare per formare lo studio di fattibilità.

        \subsection{Analisi dei documenti da produrre per la RR}
        Per avere un'idea generale sui documenti da produrre per la Revisione dei Requisiti,
        Samuele e Manuel hanno presentato alcune ricerche al gruppo.

        \subsection{Scelta dei ruoli}
        Il gruppo ha parlato di come suddividere i ruoli ai vari componenti per iniziare al meglio il progetto. Il responsabile poi si occuperà di assegnare i nuovi ruoli.
\end{document}
